\begin{definition}[正向电子的能量归一化波包]
固定传播方向和自旋，令 $p>0$，$\langle p|p'\rangle=\delta(p-p')$，
$E(p)=\sqrt{m_{\rm e}^2c^4+c^2p^2}$、$v(E)=\dd E/\dd p>0$。
能量归一化广义态和能量包络为
\begin{equation}
|E\rangle=\frac{|p(E)\rangle}{\sqrt{v(E)}},\qquad
b(E)=\frac{a(p(E))}{\sqrt{v(E)}},\qquad
\int_{m_{\rm e}c^2}^\infty|b(E)|^2\dd E=1.
\label{eq:feqo-energy-envelope-normalization}
\end{equation}
这里 $a\in L^2((0,\infty),\dd p)$，并取 $\int|a|^2\dd p=1$。
\end{definition}
\begin{proof}
由 $\delta(p(E)-p(E'))=v(E)\delta(E-E')$ 得
$\langle E|E'\rangle=\delta(E-E')$。
在 $\int a(p)|p\rangle\dd p$ 中用 $\dd p=\dd E/v(E)$ 换元，
得到 $\int b(E)|E\rangle\dd E$；对模方积分也得到相同归一化。
因此 $v^{-1/2}$ 同时由态的 delta 归一化和概率测度的换元固定。
\end{proof}

\begin{proposition}[能量平移波包之间的重叠]
若在远离谱底的能窗中，初始能量包络近似为
$b_0(E)=(2\pi\sigma_E^2)^{-1/4}\exp[-(E-E_0)^2/(4\sigma_E^2)]$，
令 $b_\ell(E)=b_0(E-\ell\hbar\omega)$，则在可忽略谱底尾概率的模型内
\begin{equation}
\langle b_\ell,b_m\rangle
=\exp\!\left[-\frac{(\ell-m)^2(\hbar\omega)^2}{8\sigma_E^2}\right].
\label{eq:feqo-sideband-wavepacket-overlap}
\end{equation}
\end{proposition}
\begin{proof}
把两个高斯指数相加并配方，公共积分中心为
$E_0+(\ell+m)\hbar\omega/2$，中心差留出上述常数指数。
公共高斯积分的归一化为一。因此 $\sigma_E\ll\hbar\omega$ 才使不同平移波包近似正交。
若精确电子谱底与某个平移包络接近，则要限制积分到 $E\ge m_{\rm e}c^2$，
同时重新计算归一化；双边平移此时不再是物理态空间上的幺正操作。
\end{proof}

真实调制作用于连续包络时，出射振幅形如
$b_{\rm out}(E)=\sum_\ell c_\ell b_0(E-\ell\hbar\omega)$。
能谱是这个总振幅的模方，包括不同 $\ell$ 的交叉项。
把每一条平移包络直接称为一个正交能量格态，需要先检查以上重叠；
仪器分辨率只能改变读出，不能把尚有重叠的输入波包变成正交态。

\begin{exbox}[单步反冲很小，传播相位仍需单独控制]
\label{chap:feqo-phenomena}
单向末态能量为 $E_0+\ell\hbar\omega$ 时，精确动量是
$p_\ell=c^{-1}\sqrt{(E_0+\ell\hbar\omega)^2-m_{\rm e}^2c^4}$。
利用反函数导数 $p'(E)=1/v$、$p''(E)=-1/(\gamma^3m_{\rm e}v^3)$，得到
\begin{equation}
p_\ell=p_0+\frac{\ell\hbar\omega}{v_0}
-\frac{(\ell\hbar\omega)^2}{2\gamma_0^3m_{\rm e}v_0^3}
+O((\ell\hbar\omega)^3).
\label{eq:feqo-sideband-momentum-recoil}
\end{equation}
取 $K_0=\SI{200}{\kilo\electronvolt}$、$\hbar\omega=\SI{2}{\electronvolt}$，
$v_0=0.6953c$，单步动量为 $1.54\times10^{-27}\,\si{\kilogram\metre\per\second}$，
约为 $p_0$ 的 $5.8\times10^{-6}$。
二阶动量项在漂移长度 $L$ 上却积累相位
$-L(\ell\hbar\omega)^2/(2\hbar\gamma_0^3m_{\rm e}v_0^3)$。
它同时随 $L$ 和 $\ell^2$ 增长；单步动量很小不能代替对最外侧边带传播相位的检查。
\end{exbox}
